\section{Introduction}
\label{sec:introduction}

Large Language Models (LLMs) have become a primary interface for answering
complex questions~\citep{openai2023gpt4,anthropic2024claude}. The interaction
format, however, remains largely linear: the model writes a long answer, and
the user scrolls through it. This format works for short factual responses, but
it becomes inefficient when the answer has internal structure. When a user asks
``compare REST and GraphQL'' or ``explain the RAG retrieval pipeline,'' they
often want to inspect one module, one tradeoff, or one stage. Plain text can
contain that information, but it provides little support for navigation.

Parallel advances in text-to-image generation~\citep{ramesh2022dalle2,saharia2022imagen,podell2024sdxl}
and vision-language understanding~\citep{liu2023llava,chen2024internvl} suggest
a better interface: \emph{render the answer as an image}. A generated
infographic, map, or scene can show modules, hierarchy, routes, and spatial
relationships at a glance. Static images, however, discard much of the
explanatory structure behind each region and provide no local follow-up
mechanism. Directly overlaying planned boxes is also fragile because image
generators do not reliably place content at the requested coordinates~\citep{feng2023layoutgpt,cheng2024intelligent}.
If the interaction layer is copied from the pre-generation plan, the user may
click an empty area or the wrong object.

\begin{figure}[t!]
    \centering
    \includegraphics[width=0.95\textwidth]{demo1.pdf}
    \caption{\textbf{\method{} turns a long-form LLM answer into an
    interactive visual image.} Given a user question, \method{} generates an
    image whose regions are clickable hotspots. Each hotspot opens a detail
    panel with the region's explanation and an in-context follow-up thread,
    so the user explores only what they need.}
    \label{fig:teaser}
    \vspace{-0.6em}
\end{figure}

We propose \method{}, a system that turns a long-form LLM answer into an
\textbf{interactive visual image}: a generated picture with transparent
hotspots over the visible regions (Figure~\ref{fig:teaser}). Each hotspot opens
the explanation for that region and carries its own follow-up thread. The
central engineering problem is alignment: a hotspot should cover the object,
panel, landmark, or diagram region the user can actually see, not merely the
box that was planned before generation.

\method{} handles this with a \textbf{generate-then-ground} pipeline. The
system turns the answer into a small set of visual modules, plans an
approximate layout, writes an image prompt, and generates the image. It then
sends the rendered image back through visual grounding. The grounding step finds where
each module actually appeared, and the frontend places the clickable hotspots
on those discovered regions. This decouples \emph{what the answer should
contain} from \emph{where the image generator rendered it}. The resulting
artifact supports region-level exploration: the user clicks a visible region,
sees the relevant detail and a local preview, and asks a follow-up without
re-explaining the entire answer. The system is provider-agnostic: it runs in a
no-key mock mode for development and demos, or proxies real text, image, and
vision models through a local backend so API keys never reach the browser.

To our knowledge, \method{} is the \textbf{first} system to convert an LLM
answer into a navigable, region-interactive visual image with vision-grounded
hotspots. Our contributions are:
\begin{itemize}
    \item \textbf{Interactive visual answers.} We introduce the task of
    rendering a long-form LLM answer as an explorable image with clickable
    hotspots, detail panels, and per-region follow-up threads.
    \item \textbf{Generate-then-ground alignment.} We decouple structured
    content from visual placement: the system first renders the image, then
    uses visual grounding and optional SAM-family mask refinement to place
    hotspots on actual visible regions rather than a hard-coded grid.
    \item \textbf{Open-source reference system.} We release a complete,
    provider-agnostic implementation with deterministic mock mode, real API
    mode, local SQLite persistence, per-hotspot follow-up threads, and a
    calibration mode for inspecting and correcting hotspot bounds.
    \item \textbf{Benchmark and grounding evaluation.} We assemble a 30-question
    benchmark spanning multiple content types and evaluate the real-provider
    pipeline with automatic completion checks and a strict visual-alignment
    gate.
\end{itemize}
